\label{sec:README}
\section{Emergent Communication in Multi-Agent Reinforcement Learning with Assembly Line Environments}

Built for a Cambridge Computer Science undergraduate dissertation project.

This open-source program provides a framework for evaluation of emergent protocols in Multi-Agent Systems.

\section{Getting Started}

Before starting, initialise the virtual environment and build the simulation environment:

\begin{verbatim}
python -m venv .venv
.venv\Scripts\activate
pip install -e .
\end{verbatim}

A set of eight configuration files are provided for customisability.

\section{Running the program}

The AI Mother Tongue communication protocol (AIM) requires pretraining a language. To train the language, run:

\begin{verbatim}
python -m controller.marl.main --mode train-language
\end{verbatim}

Once the configuration files have been customised, training starts via:

\begin{verbatim}
python -m controller.marl.main --mode train
\end{verbatim}

MARL often requires thousands or even millions of training steps to converge to an optimal policy. To accelerate this, imitation learning from an expert agent is used to pretrain the actor and critic. To use this, run:

\begin{verbatim}
python -m controller.marl.main --mode imitate
\end{verbatim}

\section{Customising Protocols}

New communication protocols can be added to the folder \texttt{controller/marl/comms/}. The \texttt{comms.py} file contains an abstract class that new protocols should inherit from.

\section{Analysis}

A large set of evaluation tools are provided to analyse communication between agents. This includes:

\begin{itemize}[noitemsep, topsep=0pt]
    \item Discrete codebook usage
    \item Discrete codebook PCA
    \item Comparing training from multiple runs
    \item Memory probing
    \item Conditional probability of target or action given code
\end{itemize}
